\section{Conclusión individual: Washington Apunte}

Mi parte fue el contrato \texttt{servicio.proto} y el servidor gRPC. En el
contrato definí el servicio \texttt{ComunicacionNodos}, con la función
\texttt{EnviarMensaje}, y los mensajes \texttt{Mensaje} y \texttt{Acuse}. En el
servidor implementé esa función: valida la solicitud, actualiza el reloj de
Lamport, registra los eventos \texttt{receive} y \texttt{send} y devuelve el
acuse. Las llamadas se atienden con un grupo fijo de 10 hilos.

La decisión que más tuve que justificar fue usar \texttt{int64} para los tres
timestamps, aunque el reloj nunca es negativo. Lo elegí porque facilita la
compatibilidad con lenguajes sin enteros sin signo nativos y porque permite que
el servidor reciba un valor negativo y lo rechace con
\texttt{INVALID\_ARGUMENT}, como comprueba \texttt{test\_timestamp\_negativo}.

Escribí 6 pruebas para el servidor, que forman parte de las 56 del proyecto.
La de concurrencia lanza 30 llamadas desde 10 hilos y ningún timestamp se
repite, pero aprendí a no exagerar ese resultado: demuestra que los valores son
únicos, no que cada llamada sea atómica. También vi que el contrato fija los
tipos, pero no los valores: el emisor vacío y el timestamp negativo los sigue
validando mi servidor.
